\chapter{Граничный родительский след: совместимость секторов Вильсона, дефекта и поколений}

После отрицательных результатов для носителя и минимального треугольника
редукций проверяется граничный источник. Вопрос состоит в том, существуют ли
одно минимальное граничное гильбертово пространство, одна сохраняющая
симметрии алгебра, один след или суперслед и один квадратичный родительский
оператор, из которых одновременно следуют вильсоновский сектор, конденсат
спаривания, тетраэдрическая ось поколений и ровно одна мода Майораны.

\section{Зафиксированные требования}

До построения общего носителя требуются: канонический источник нечётной
ветви Вильсона; ненулевой конденсат заряда два; голономия \(2\pi/3\) и ось
ориентированного трёхцикла; ровно одна вещественная нулевая мода Майораны
после связи поколений; второй независимый результат, чувствительный к
нормировке. Коэффициент, топологическая метка, контур или проектор,
вставленные после сравнения, означают неудачу.

\section{Сильнейшие готовые блоки}

Шестнадцатимерный гауссов блок Вильсона точно воспроизводит нелинейную пару
коэффициентов
\begin{equation}
 -\frac13\log Z_J
 =\frac8{1-\cos\theta}+\frac83\log(1-\cos\theta)+\text{пост.},
\end{equation}
но требует проектора фиксированного заряда либо мнимого когерентного
источника, зависящего от оси.

Проекторная суперкривизна выбирает четыре тетраэдрических проектора, а
намотка \(\nu=\pm1\) превращает их в восемь ориентированных трёхциклов с
голономией \(2\pi/3\). Нормированное отображение момента колчана даёт
устойчивый конденсат, однако его вложение KO6 порождает
\((\rho^2+r^2)^2\), а не \((\rho^2-r^2)^2\). Поколенческий генератор
\(K_n=(2\pi/3)[n]_\times\) уменьшает три слепые к поколениям моды ядра до
одной, поскольку \(\dim\ker K_n=1\), но корневой источник, конденсат и
ограничение связи остаются условными.

\section{Развилка одного следа}

Пусть сохраняющая симметрии граничная алгебра имеет разложение
\begin{equation}
 \mathcal A_\partial=\bigoplus_{a=1}^{k}M_{n_a}(\mathbb F_a).
\end{equation}
Всякий положительный след равен
\begin{equation}
 \tau_\partial=\sum_{a=1}^{k}w_a\Tr_a,
\end{equation}
и после условия \(\tau_\partial(1)=1\) остаются \(k-1\) относительных весов.
Текущие вильсоновский, колчанный KO6 и ядерный модули неэквивалентны по
заряду, степени формы, градуировке и вещественности. Поэтому их прямая сумма
имеет не менее трёх центральных блоков и двух свободных весов.

Замена прямой суммы полным матричным фактором дала бы один след, но ввела бы
новые внедиагональные связующие поля между неэквивалентными секторами и новое
действие.

\section{Препятствие фиксированного заряда}

Канонический проектор заряда имеет вид
\begin{equation}
 P_q=\frac1{2\pi}\int_0^{2\pi}d\alpha\,
 e^{i\alpha(\widehat Q-q)},
 \qquad Z_q=\Tr(P_qe^{-\beta H}).
\end{equation}
Одна диагональная связь Гаусса на восьми импульсах имеет ранг один и
оставляет семь направлений; две блочные связи оставляют шесть. Требуемый
вектор \((1,1,1,3,3,3,3,3)\) однозначно фиксируется только ограничением
ранга восемь либо проектором, уже содержащим этот вектор.

\begin{theorem}[След не выбирает сектор заряда]
Один след и одна диагональная связь Гаусса не выводят требуемую
восьмикомпонентную нечётную ветвь. Проектор является дополнительными
дискретными данными, если он не следует из отсутствующей структуры связей
высшего ранга.
\end{theorem}

\section{Круг в сильнейшем объединении}

Имеющиеся модули образуют логическую цепь
\(n\to J(n)\to V_W(\theta,n)\), но сама ось \(n\) выбирается другой
проекторной кривизной. Прямое сложение проекторного, гауссова, колчанного и
ядерного блоков не создаёт общего механизма: их относительные нормировки и
связующие члены остаются независимыми.

Формализм BV--BFV и теория детерминантных линий поддерживают язык граничной
склейки (\path{arXiv:2012.13983}, \path{arXiv:hep-th/9405012}), но не
выбирают специфические кратности, проектор заряда, потенциал спаривания и
скалярную тривиализацию проекта.

\section{Вывод}

\begin{theorem}[Результат проверки граничного следа]
Ни одно построенное граничное пространство с сохраняющим симметрии следом не
выводит одновременно все пять требований. Каждый отдельный результат точен
или условен локально, но общей меры и общего правила выбора нет.
\end{theorem}

\begin{protocol}[Двоичная фиксация]
\begin{enumerate}
 \item общая граничная кинематика: \textbf{частично выполнена};
 \item точная нелинейная пара Вильсона: \textbf{получена на уровне оператора};
 \item один родительский след: \textbf{не получен};
 \item выбор заряда или источника: \textbf{не выведен};
 \item совместный вывод конденсата, оси и моды Майораны:
 \textbf{не выполнен};
 \item математическое и физическое замыкание: \textbf{не достигнуты}.
\end{enumerate}
\end{protocol}

Текущая граничная реализация закрыта как родительская архитектура версии V.
Следующая проверка --- \path{version5_finite_geometry_complexity_bound_gate}.
Машинный реестр сохранён в
\path{s2t_v5_boundary_parent_trace_freeze_gate_results.json}.